% =============================================================
% Epilogue: Summary, Relativistic Corrections, and Open Problems
% =============================================================

\chapter{Epilogue: Relativistic Corrections to Classical Instabilities}
\label{ch:epilogue}

\section{Summary of Relativistic Modifications}

The preceding fourteen chapters have systematically extended every classical
hydrodynamic and hydromagnetic instability treated by Chandrasekhar into the
regime of special- and general-relativistic fluids, using the causal,
first-order BDNK viscous hydrodynamics framework throughout.  The central
lesson is that \emph{every} classical instability criterion acquires
relativistic corrections that depend on the ratio of the fluid's enthalpy
density to its rest-mass energy density, $\enthalpy / \rdensity c^2$, and on
the flow Lorentz factor $\Lf$.  These corrections are not merely kinematic
rescalings: the BDNK causal structure introduces new time-scales
($\tauq$, $\taupi$) that modify overstability boundaries and can suppress or
enhance instabilities that are absent in the Navier--Stokes limit.

Below we collect the principal results in tabular form.

\section{Master Table of Relativistic Correction Factors}
\label{sec:master-table}

Table~\ref{tab:master-corrections} summarises, for each instability, the
leading-order relativistic correction factor that multiplies the classical
threshold or growth rate.  In every case the non-relativistic limit
$v/c \to 0$ recovers Chandrasekhar's original result.

\begin{table}[htbp]
\centering
\small
\caption{Summary of relativistic corrections to classical instabilities.
Here $\Lf$ is the Lorentz factor, $\enthalpy = \edensity + p$ the relativistic
enthalpy density, $\rdensity$ the rest-mass density, $\cs$ the relativistic
sound speed, and $\tauq$, $\taupi$ the BDNK relaxation time-scales.}
\label{tab:master-corrections}
\renewcommand{\arraystretch}{1.35}
\begin{tabular}{@{}p{3.8cm}cp{5.5cm}@{}}
\toprule
\textbf{Instability (Chapter)} & \textbf{Key correction factor} & \textbf{Physical origin} \\
\midrule
Thermal convection (II--III)
  & $\Rarel = \Ra\,\dfrac{\rdensity}{\enthalpy/c^2}$
  & Inertia of enthalpy replaces rest mass \\[6pt]
Thermal convection with rotation (IV--V)
  & $\Tarel = \Ta\,\dfrac{\rdensity^2}{(\enthalpy/c^2)^2}$
  & Relativistic Coriolis coupling \\[6pt]
Thermal convection with magnetic field (VI)
  & $\Qrel = Q\,\dfrac{\rdensity}{\enthalpy/c^2}\,\dfrac{1}{1+\bsq/\enthalpy}$
  & Magnetic enthalpy contribution \\[6pt]
Rayleigh--Taylor (X)
  & $\Lf^2\,\dfrac{\enthalpy}{\rdensity c^2}$
  & Relativistic effective inertia \\[6pt]
Kelvin--Helmholtz (XI)
  & $\Lf_1^2\Lf_2^2\,\dfrac{\enthalpy_1\,\enthalpy_2}{\rdensity_1\rdensity_2 c^4}$
  & Lorentz-boosted interface momentum flux \\[6pt]
Stability of Couette flow (VII)
  & $\dfrac{\Lf^2\,\enthalpy}{\rdensity c^2}$
  & Relativistic angular momentum density \\[6pt]
Stability of viscous flow (VIII--IX)
  & BDNK: $1 + \taupi\,\omega$
  & Causal viscous propagation delay \\[6pt]
Jeans / gravitational (XIII)
  & $k_J^2 \to k_J^2\,\dfrac{\enthalpy}{\rdensity c^2}\bigl(1+\frac{3p}{\edensity}\bigr)$
  & Pressure gravitates (Tolman--Oppenheimer--Volkoff) \\[6pt]
Relativistic jets / current-driven (XII)
  & $\Lf^2\bigl(1 + \vA^2/c^2\bigr)$
  & Magnetic tension boosted by flow Lorentz factor \\[6pt]
Variational / energy principles (XIV)
  & $\delta^2 W \to \delta^2 W + \order{v^2/c^2}$
  & Relativistic energy functional corrections \\
\bottomrule
\end{tabular}
\end{table}

\section{Open Problems and Future Directions}
\label{sec:open-problems}

\begin{enumerate}
\item \textbf{Nonlinear saturation in the relativistic regime.}
  The present work is entirely linear.  Nonlinear saturation amplitudes for
  relativistic Rayleigh--Taylor, Kelvin--Helmholtz, and magneto-rotational
  instabilities remain poorly constrained analytically.

\item \textbf{BDNK transport coefficients from kinetic theory.}
  The BDNK frame coefficients $(\epsone, \betaone, \alphaone)$ used here are
  treated as free parameters subject to causality and stability constraints.
  First-principles derivations from relativistic Boltzmann theory for
  astrophysical plasmas (as opposed to quark--gluon plasmas) are still needed.

\item \textbf{General-relativistic extensions in strong-field backgrounds.}
  We have worked in Minkowski or weakly curved backgrounds.  Extending these
  results to Kerr spacetimes---relevant for accretion disks around spinning
  black holes---requires coupling the BDNK equations to the full Einstein
  field equations.

\item \textbf{Numerical verification via GRMHD simulations.}
  The analytic growth rates and threshold corrections derived here provide
  quantitative predictions testable in state-of-the-art general-relativistic
  magnetohydrodynamic codes such as \texttt{BHAC}, \texttt{IllinoisGRMHD},
  \texttt{AthenaK}, and \texttt{SpECTRE}.  Systematic comparison between
  linear theory and nonlinear simulation is a high-priority next step.

\item \textbf{Instabilities in neutron-star mergers.}
  The Kelvin--Helmholtz and magneto-rotational instabilities at the
  merger interface of binary neutron stars operate at Lorentz factors
  $\Lf \sim 1$--$2$ and temperatures $T \sim 50$\,MeV.  The BDNK corrections
  derived in Chapters~X--XI and VII are directly applicable and could refine
  sub-grid models used in numerical relativity simulations.

\item \textbf{Resistive and radiative extensions.}
  This monograph treats ideal MHD with viscous and thermal dissipation.
  Resistive and optically thick (radiation-MHD) extensions of the BDNK
  framework---and their implications for instability thresholds---remain
  largely unexplored.

\item \textbf{Turbulent transport in relativistic convection.}
  The relativistic Rayleigh and B\'enard problems derived here set the stage
  for a relativistic mixing-length theory.  Connecting the linear onset
  criteria to turbulent transport coefficients in relativistic stellar
  envelopes is an important astrophysical application.
\end{enumerate}

\section{Connection to Numerical GRMHD Simulations}
\label{sec:grmhd-connection}

Modern GRMHD codes solve the full nonlinear equations of relativistic
magnetohydrodynamics on adaptive meshes, often coupled to nuclear
equations of state and neutrino transport.  The linear-stability results of
this book connect to simulations in three ways:

\paragraph{Benchmark tests.}
The analytic dispersion relations derived here---particularly for the
relativistic Rayleigh--Taylor (Chapter~X), Kelvin--Helmholtz (Chapter~XI),
and magnetorotational (Chapter~VII) instabilities---provide exact or
semi-analytic growth rates that can serve as convergence tests for GRMHD
codes.  A simulation that cannot reproduce the correct linear growth rate in
the appropriate limit has a bug or insufficient resolution.

\paragraph{Sub-grid models.}
In large-scale astrophysical simulations (e.g., jet propagation, accretion
disk evolution), the smallest unstable scales are unresolved.  The
relativistic correction factors tabulated in
Table~\ref{tab:master-corrections} can inform sub-grid prescriptions for
turbulent viscosity and mixing, replacing \emph{ad hoc} Newtonian estimates
with physically motivated relativistic expressions.

\paragraph{Interpreting simulation results.}
When a simulation exhibits an instability---for instance, kink modes in a
relativistic jet---the analytic framework of Chapter~XII allows one to
identify the dominant mode, predict its growth rate, and assess whether the
observed behaviour is physical or numerical.

\bigskip
\noindent
In summary, this monograph provides the theoretical scaffolding for
understanding hydrodynamic and hydromagnetic instabilities in relativistic
astrophysical flows.  The BDNK formalism ensures that the dissipative sector
is causal and stable, resolving long-standing pathologies of relativistic
Navier--Stokes theory.  The systematic chapter-by-chapter correspondence with
Chandrasekhar's classical treatise is intended to make the relativistic
generalisations as accessible as possible to researchers trained in the
Newtonian theory.
